% Partie : histoire
% Chapitre : biographies
% Section : Julia Robinson

\documentclass[../../../include/open-logic-section]{subfiles}

\begin{document}

\olfileid[fr]{his}{bio}{rob}

\olsection{Julia Robinson}

\olphoto{robinson-julia}{Julia Robinson}

Julia Bowman Robinson était une mathématicienne américaine.
Elle est surtout connue pour ses travaux sur les problèmes de décision,
et notamment pour sa contribution à la résolution du dixième problème
de Hilbert. Elle naquit à Saint-Louis, dans le Missouri, le
8 décembre 1919. Elle se souvenait d'avoir été fascinée par les
nombres dès son enfance \citep[4]{Reid1986}. À neuf ans, elle contracta
la scarlatine et subit plusieurs rechutes de rhumatisme articulaire
aigu. Elle dut passer beaucoup de temps alitée et prit du retard dans
ses études. Des précepteurs l'aidèrent à le rattraper, mais les
séquelles physiques de sa maladie marquèrent durablement sa vie.

Malgré ces épreuves, Robinson termina ses études secondaires avec
plusieurs prix en mathématiques et en sciences. Elle commença ses
études supérieures au San Diego State College, puis passa sa dernière
année de premier cycle à l'université de Californie à Berkeley. Elle
y fit la connaissance du mathématicien Raphael Robinson. Ils devinrent
de bons amis et se marièrent en 1941. Parce qu'elle était l'épouse
d'un membre du corps enseignant, Julia Robinson n'avait pas le droit
d'enseigner au département de mathématiques de Berkeley. Elle continua
à suivre des cours de mathématiques en auditrice libre, tout en
espérant quitter l'université pour fonder une famille. Peu après son
mariage, cependant, elle contracta une pneumonie. Un médecin lui
apprit que le rhumatisme articulaire aigu de son enfance avait laissé
d'importantes cicatrices sur son cœur. Il prédit qu'elle ne vivrait
pas au-delà de quarante ans et lui déconseilla d'avoir des enfants
\citep[13]{Reid1986}.

Robinson traversa une longue période de dépression, puis décida de
reprendre ses études de mathématiques. De retour à Berkeley, elle
obtint son doctorat en 1948 sous la direction d'Alfred Tarski. Tarski
avait montré que la théorie du premier ordre des nombres réels était
décidable, tandis qu'il découlait des travaux de Gödel que celle des
entiers naturels ne l'était pas. La décidabilité de la théorie du
premier ordre du corps des nombres rationnels demeurait une grande
question ouverte. Robinson en établit l'indécidabilité dans sa thèse.
Le résultat fut publié l'année suivante \citeyearpar{Robinson1949}.

S'intéressant aux problèmes de décision, Robinson s'attaqua ensuite
au dixième problème de Hilbert, l'un des 23 célèbres problèmes posés
par David Hilbert en 1900. Il demande s'il existe un algorithme
qui, pour toute équation polynomiale à coefficients entiers, décide
en un temps fini si elle admet une solution en nombres entiers.
On pourrait par exemple lui soumettre l'équation $3x^2 - 2y +3 = 0$.
De telles questions sont appelées
\emph{problèmes diophantiens}. Après quelques premiers succès,
Robinson s'associa à Martin Davis et Hilary Putnam, qui travaillaient
eux aussi sur ce problème. Ils montrèrent que les problèmes
diophantiens exponentiels, où les inconnues peuvent également
apparaître en exposant, sont indécidables. Ils montrèrent aussi
que l'existence d'une relation diophantienne à croissance
exponentielle, connue sous le nom d'hypothèse de Julia Robinson
(«~J.R.~»), impliquait l'indécidabilité du dixième problème de Hilbert
\citep{DavisPutnamRobinson1961}. Robinson poursuivit ses recherches
sur ce problème tout au long des années 1960. En 1970, le jeune
mathématicien russe Yuri Matijasevich démontra enfin la conjecture
J.R. Le résultat obtenu est aujourd'hui appelé théorème de
Matijasevich--Robinson--Davis--Putnam, ou, en abrégé, théorème MRDP.
Matijasevich et Robinson devinrent amis et cosignèrent plusieurs
articles. Dans une lettre à Matijasevich, Robinson écrivit un jour :
«~En fait, je suis très heureuse que notre travail commun (à des
milliers de kilomètres l'un de l'autre) nous fasse manifestement
avancer davantage que chacun de nous n'y parviendrait seul~»
\citep[45]{Matijasevich1992}.

Robinson fut la première femme à présider l'American Mathematical
Society et la première femme élue à la section de mathématiques de la National Academy of Sciences.
Elle mourut le 30 juillet 1985, à 65 ans, après qu'on lui eut
diagnostiqué une leucémie.

\begin{reading}
Les articles mathématiques de Robinson sont réunis dans ses
\textit{Collected Works} \citep{Robinson1996}, qui comprennent aussi
une réimpression de sa notice biographique de la National Academy of
Sciences \citep{Feferman1994}. Sa sœur aînée, Constance Reid, publia
une «~Autobiography of Julia~» fondée sur des entretiens
\citep{Reid1986}, ainsi qu'une biographie complète
\citep{Reid1996}. George Csicsery réalisa un court documentaire sur
Robinson et le dixième problème de Hilbert \citep{Csicsery2016}.
Pour un bref témoignage sur les collaborations de Yuri Matijasevich
avec Robinson et sur l'influence de celle-ci sur ses travaux, voir
\citep{Matijasevich1992}.
\end{reading}

\end{document}
